\documentclass[11pt]{article}
\usepackage[margin=1in]{geometry}
\usepackage{amsmath,amssymb,amsthm}
\usepackage{booktabs}
\usepackage{hyperref}

\newtheorem{theorem}{Theorem}
\newtheorem{lemma}[theorem]{Lemma}
\newtheorem{proposition}[theorem]{Proposition}
\theoremstyle{definition}
\newtheorem{definition}[theorem]{Definition}
\theoremstyle{remark}
\newtheorem{remark}[theorem]{Remark}

\title{Disproof of Conjecture 00000000603:\\
the Conway polynomial of $T(2,25)$ does not peak at $\lfloor (p-1)(q-1)/4\rfloor$}
\author{jilint777}
\date{2026-10-04}

\begin{document}
\maketitle

\begin{abstract}
Conjecture 00000000603 asserts that for every torus knot $T(p,q)$ the
coefficient sequence of the Conway polynomial $\nabla(z)$ is nonnegative,
has no internal zeros, is log-unimodal, and has its peak at position
$\lfloor (p-1)(q-1)/4\rfloor$. We show that the peak-position clause is false.
For $T(2,25)$ the coefficients of $z^0,z^2,\dots,z^{24}$ are
$1,78,1001,5005,12870,19448,18564,\dots$, so the peak is at index $5$
(the term $z^{10}$), while $\lfloor 1\cdot 24/4\rfloor=6$. For $T(2,61)$
the predicted position $15$ fails under every reasonable indexing convention.
Everything is verified in a self-contained Lean~4 project and an independent
Python script.
\end{abstract}

\section{The conjecture}

The statement (English version) reads:
\begin{quote}
\emph{Definition:} The specialization $P_K(1,z)$ of the HOMFLY polynomial
$P_K(a,z)$ is the Conway polynomial $\nabla_K(z)$.
\emph{Conjecture:} For the torus knot $T(p,q)$, the coefficient sequence of
$\nabla(z)$ is nonnegative, has no internal zeros, and is unimodal on the
logarithmic scale; the peak position is $\lfloor (p-1)(q-1)/4\rfloor$.
\end{quote}
The Chinese version says the same thing. The statement is a conjunction about
every torus knot, so a single torus knot where the peak is not at
$\lfloor (p-1)(q-1)/4\rfloor$ refutes it.

\paragraph{Readings of ``peak position''.}
The Conway polynomial of a knot contains only even powers of $z$, and
$\deg\nabla_{T(p,q)} = 2g = (p-1)(q-1)$. So $\lfloor (p-1)(q-1)/4\rfloor$ is
naturally an index into the sequence of even coefficients. We consider three
readings of ``the coefficient at position $i$'':
\begin{enumerate}
\item[(R1)] the coefficient of $z^{i}$;
\item[(R2)] the coefficient of $z^{2i}$ (even coefficients, indexed from $0$);
\item[(R3)] the coefficient of $z^{2(i-1)}$ (even coefficients, indexed from $1$).
\end{enumerate}
Reading (R2) is the natural one: it is the only one under which the
formula holds for all of $T(2,3),\dots,T(2,19)$ (Table~\ref{tab:small}).
We use the weakest meaning of ``the peak is at $i$'': the coefficient at
position $i$ is \emph{a} maximum, with ties allowed. We show that a strictly
larger coefficient exists.

\section{\texorpdfstring{The Conway polynomial of $T(2,n)$}{The Conway polynomial of T(2,n)}}

Let $L_n$ be the closure of the 2-strand braid $\sigma_1^n$. For odd $n$,
$L_n$ is the torus knot $T(2,n)$. Also $L_0$ is the two-component unlink and
$L_1$ is the unknot. The Conway polynomial satisfies
$\nabla(\text{unknot})=1$, $\nabla(\text{split link})=0$, and the skein relation
$\nabla(L_+)-\nabla(L_-)=z\,\nabla(L_0)$. Apply the skein relation at one
crossing of $\sigma_1^{n+2}$. Changing that crossing gives $\sigma_1^{n}$, and
smoothing it gives $\sigma_1^{n+1}$. Hence
\begin{equation}\label{eq:skein}
F_0=0,\qquad F_1=1,\qquad F_{n+2}=F_n+z\,F_{n+1},\qquad F_n:=\nabla(L_n).
\end{equation}
These conditions determine every $F_n$. The $F_n$ are the Fibonacci polynomials.

\begin{lemma}\label{lem:closed}
$F_n(z)=\sum_{j\ge 0}\binom{n-1-j}{j} z^{\,n-1-2j}$ for $n\ge1$. In particular,
for $n=2k+1$,
\[
\nabla_{T(2,2k+1)}(z)=\sum_{i=0}^{k}\binom{k+i}{2i}\,z^{2i}.
\]
\end{lemma}
\begin{proof}
Write $G_n$ for the right-hand side. Then $G_1=1$ and $G_2=z$, and
\eqref{eq:skein} gives $F_2=F_0+zF_1=z$. For $n\ge1$, the coefficient of
$z^{n+1-2j}$ in $G_n+zG_{n+1}$ is
$\binom{n-j}{j-1}+\binom{n-j}{j}=\binom{n+1-j}{j}$ by Pascal's rule, which is the
coefficient in $G_{n+2}$. So $G_n$ satisfies the same recursion as $F_n$,
and $F_n=G_n$ by induction. For $n=2k+1$, put $i=k-j$. The exponent is
$2k-2j=2i$ and the coefficient is $\binom{2k-j}{j}=\binom{k+i}{k-i}=\binom{k+i}{2i}$.
\end{proof}

As an independent check, the Alexander polynomial of a torus knot is
$\Delta_{T(p,q)}(t)=\frac{(t^{pq}-1)(t-1)}{(t^p-1)(t^q-1)}$, and
$\nabla(t^{1/2}-t^{-1/2})=\Delta(t)$. The script \texttt{verify.py} derives the
Conway polynomial from this formula and confirms that it agrees with
\eqref{eq:skein} and Lemma~\ref{lem:closed} for all $T(2,n)$, $n\le 61$.

\begin{table}[ht]
\centering
\begin{tabular}{rlcc}
\toprule
$n$ & coefficients of $z^0,z^2,z^4,\dots$ in $\nabla_{T(2,n)}$ & peak index & $\lfloor (n-1)/4\rfloor$\\
\midrule
3 & 1, 1 & 0, 1 & 0\\
5 & 1, 3, 1 & 1 & 1\\
7 & 1, 6, 5, 1 & 1 & 1\\
9 & 1, 10, 15, 7, 1 & 2 & 2\\
13 & 1, 21, 70, 84, 45, 11, 1 & 3 & 3\\
\textbf{25} & 1, 78, 1001, 5005, 12870, \textbf{19448}, 18564, 11628, \dots & \textbf{5} & \textbf{6}\\
\bottomrule
\end{tabular}
\caption{Even Conway coefficients of $T(2,n)$ and the predicted peak index.}
\label{tab:small}
\end{table}

\section{The counterexamples}

\begin{theorem}\label{thm:25}
For $T(2,25)$ we have $(p-1)(q-1)/4 = 6$, and
\begin{align*}
\nabla_{T(2,25)}(z)={}&1+78z^2+1001z^4+5005z^6+12870z^8+19448z^{10}+18564z^{12}\\
&+11628z^{14}+4845z^{16}+1330z^{18}+231z^{20}+23z^{22}+z^{24}.
\end{align*}
The coefficient $19448$ of $z^{10}$ is strictly larger than the coefficient
$18564$ of $z^{12}$ (position $6$ under (R2)) and the coefficient $5005$ of $z^6$
(position $6$ under (R1)). So the peak-position clause fails for $T(2,25)$
under readings (R1) and (R2).
\end{theorem}
\begin{proof}
By Lemma~\ref{lem:closed} with $k=12$, the coefficient of $z^{2i}$ is
$\binom{12+i}{2i}$. This gives the list above. For example,
$\binom{17}{10}=19448$, $\binom{18}{12}=18564$ and $\binom{15}{6}=5005$.
\end{proof}

For (R2), the ratio of consecutive coefficients is
$\binom{k+i+1}{2i+2}\big/\binom{k+i}{2i}=\frac{(k+i+1)(k-i)}{(2i+1)(2i+2)}$.
For $k=12$ and $i=5$ it equals $18\cdot 7/(11\cdot 12)=126/132<1$. So the
sequence is already decreasing from index $5$ to index $6$. In general the
ratio is about $(k^2-i^2)/(4i^2)$, so the peak index of $\binom{k+i}{2i}$ is
about $k/\sqrt5\approx 0.447\,k$. The predicted index is
$\lfloor (n-1)/4\rfloor=\lfloor k/2\rfloor$, so the gap between them grows
linearly in $k$. $T(2,25)$ and $T(3,13)$ (both of genus 12) are the
smallest-genus torus knots that violate (R2), according to a survey of all
$T(p,q)$ with $p<q\le 30$ in \texttt{verify.py}.

\begin{theorem}\label{thm:61}
For $T(2,61)$ we have $\lfloor 60/4\rfloor=15$. The unique largest coefficient of
$\nabla_{T(2,61)}$ is $[z^{26}]=\binom{43}{26}=421171648758$, whereas
$[z^{15}]=0$, $[z^{28}]=\binom{44}{28}=416714805914$ and
$[z^{30}]=\binom{45}{30}=344867425584$. So the predicted position $15$ is not
a peak under any of (R1), (R2), (R3).
\end{theorem}
\begin{proof}
Lemma~\ref{lem:closed} with $k=30$ gives these values. The coefficient of
$z^{15}$ is zero because only even powers occur. The values are checked
exactly in Lean and in Python.
\end{proof}

\begin{remark}
Reading (R3) happens to hold for $T(2,25)$ (position 6 is $z^{10}$). It
already fails for $T(2,5)$, where position 1 is $z^0$ with coefficient
$1<3$. It also fails for $T(2,61)$. The other clauses of the conjecture
(nonnegativity, log-unimodality) are not needed: the conjecture is a
conjunction, and one false clause refutes it.
\end{remark}

\begin{theorem}
Conjecture 00000000603 is false.
\end{theorem}

\section{Lean formalization}

The directory \texttt{lean4/} contains a self-contained Lean~4.19.0 project
that uses only the core library (no Mathlib). Its contents are:
\begin{itemize}
\item \texttt{ConwaySkein c}: a family of polynomials in $z$, given as
coefficient functions $\mathbb N\to\mathbb Z$, that satisfies \eqref{eq:skein}
(unlink, unknot, skein relation for $\sigma_1^{n+2}$).
\item \texttt{skein\_unique}: every such family agrees, coefficient by
coefficient, with the reference lists \texttt{ref n} computed by the
recursion. \texttt{ref\_skein}: the reference family satisfies the
hypotheses, so they are consistent and the result is not vacuous.
\item \texttt{ref\_25}: the full coefficient list of $\nabla_{T(2,25)}$.
\texttt{ref\_61\_peak}: in $\nabla_{T(2,61)}$ the coefficient of $z^{26}$ is
strictly larger than every other coefficient. Both are proved by kernel
\texttt{decide}.
\item \texttt{PeakClaim c pos}: the conjecture's peak clause for all $T(2,n)$
($n$ odd, $n\ge3$), with position $\lfloor (2-1)(n-1)/4\rfloor$ and the
indexing convention \texttt{pos} $\in\{$(R1), (R2), (R3)$\}$.
\item \texttt{conjecture\_00000000603\_false}: for every family satisfying
the skein facts, \texttt{PeakClaim} fails under all three readings.
\texttt{T\_2\_25\_counterexample} records the smallest example for (R1) and (R2).
\end{itemize}
\texttt{lake build} succeeds. There is no \texttt{sorry}, no
\texttt{native\_decide}, and no added axioms. \texttt{\#print axioms} reports at
most \texttt{propext} and \texttt{Quot.sound}. The coefficient computations
\texttt{ref\_25} and \texttt{ref\_61\_peak} depend on no axioms at all.

\section{Reproduction}
\begin{verbatim}
cd lean4 && lake build       # Lean 4.19.0, core only
python3 verify.py            # standard library only
pdflatex report.tex
\end{verbatim}

\begin{thebibliography}{9}
\bibitem{conway} J.~H.~Conway, An enumeration of knots and links, and some of
their algebraic properties, in \emph{Computational Problems in Abstract
Algebra}, Pergamon, 1970, 329--358.
\bibitem{kauffman} L.~H.~Kauffman, The Conway polynomial,
\emph{Topology} 20 (1981), 101--108.
\bibitem{rolfsen} D.~Rolfsen, \emph{Knots and Links}, Publish or Perish, 1976
(Alexander polynomial of torus knots).
\end{thebibliography}

\end{document}
